线性规划的目标函数和约束均为线性式，变量取实数。本节统一使用
\[
\max c^Tx,\qquad Ax\le b,\quad x\ge0,
\]
其中 $A$ 有 $m$ 行、$n$ 列，分别对应约束和变量。

常见形式可作如下转换：
\begin{itemize}[leftmargin=*,nosep]
\item 最小化 $c^Tx$：改为最大化 $-c^Tx$，最优值取负。
\item $a^Tx\ge b$：两边取负，得到 $-a^Tx\le-b$。
\item $a^Tx=b$：拆成 $a^Tx\le b$ 与 $-a^Tx\le-b$。
\item 自由变量 $x_j$：令 $x_j=u_j-v_j$，其中 $u_j,v_j\ge0$。
\item $x_j\le0$：令 $x_j=-z_j$，其中 $z_j\ge0$。
\item $x_j\ge l_j$：令 $x_j=l_j+z_j$，其中 $z_j\ge0$。若所有变量都有下界，右端变为 $b-Al$，目标值加回 $c^Tl$。
\item $l_j\le x_j\le u_j$：平移后再加入 $z_j\le u_j-l_j$。
\end{itemize}
求解后按上述代换恢复原变量。目标函数中的常数不影响最优解，计算最优值时需要加回。

线性规划可能无解、具有有限最优值，或目标无界。可行域无界不代表目标无界，例如 $\max(-x)$、$x\ge0$ 的最优值为 $0$。

整数系数不保证整数最优解。例如 $\max x$、$2x\le1$、$x\ge0$ 的最优值为 $1/2$，要求 $x$ 为整数时最优值为 $0$。去掉整数限制得到线性规划松弛；最大化时，松弛的最优值给出整数最优值的上界。
